% part: intuitionistic-logic
% chapter: soundness-completeness
% section: lindenbaum

\documentclass[../../../include/open-logic-section]{subfiles}

\begin{document}

\olfileid{int}{sc}{lin}

\olsection{लिंडेनबाउमचे पूर्वप्रमेय}

अंतःप्रज्ञावादी तर्कशास्त्राचे संपूर्णता प्रमेय सिद्ध करताना
$\Gamma \Proves/ !A$ असे गृहीत धरून
$\mSat{M}{\Gamma}$ आणि~$\mSat/{M}{!A}$
असे प्रतिमान रचले जाते.

अभिजात तर्कशास्त्रात !!{derivability} चा संबंध
सुसंगततेच्या संकल्पनेत आणता येतो. कारण !!{formula}s च्या
एखाद्या संचापासून !!a{formula}~$!A$
!!{derivable} असते तेव्हा आणि केवळ तेव्हाच
त्या संचात~$!A$ चे नकरण घातल्यावर मिळणारा
संच विसंगत असतो. अंतःप्रज्ञावादी तर्कशास्त्रात
हे शक्य नाही. अंतःप्रज्ञावादी तर्कशास्त्रात
$\lnot!A$ विसंगत असल्यास केवळ
$\Proves \lnot\lnot !A$ हे मिळते.
$\lnot\lnot!A \lif !A$ हे सर्वसाधारणपणे
अंतःप्रज्ञावादी रीतीने खरे नसल्याने
$\Proves !A$ असा निष्कर्ष काढता येत नाही.

म्हणून प्रतिमान~$\mModel{M}$ रचताना
!!{formula}~$!A$ ची !!{derivability} नसणे
याची नोंद ठेवावी लागेल. त्यामुळे अभिजात
पद्धतीप्रमाणे प्रतिमान~$\mModel{M}$
रचण्यासाठी $\Gamma^* \supseteq \Gamma$
असा संपूर्ण संच वापरता येणार नाही;
कारण प्रत्येक संपूर्ण संच~$\Gamma^*$ साठी
$\Gamma^* \Proves !A \lor \lnot !A$.

संपूर्ण संच~$\Gamma^*$ वापरण्याऐवजी
सूत्रांच्या अभाज्य संचाची संकल्पना वापरू:

\begin{defn}\ollabel{defn:prime}
  !!{formula}s चा संच~$\Gamma$
  \emph{अभाज्य} असतो तेव्हा आणि केवळ तेव्हाच
  पुढील अटी पूर्ण होतात:
  \begin{enumerate}
  \item\ollabel{defn:prime1} $\Gamma$ सुसंगत असतो,
    म्हणजे $\Gamma \Proves/ \lfalse$;
  \item\ollabel{defn:prime2} $\Gamma \Proves !A$
    असेल, तर $!A \in \Gamma$; आणि
  \item\ollabel{defn:prime3} $!A \lor !B \in \Gamma$
    असेल, तर $!A \in \Gamma$ किंवा
    $!B \in \Gamma$.
  \end{enumerate}
\end{defn}

\begin{lem}[लिंडेनबाउमचे पूर्वप्रमेय]
  \ollabel{lem:lindenbaum}
  $\Gamma \Proves/ !A$ असेल, तर
  $\Gamma^* \supseteq \Gamma$ असा
  अभाज्य~$\Gamma^*$ अस्तित्वात असतो
  आणि $\Gamma^* \Proves/ !A$.
\end{lem}

\begin{proof}
  $!B \lor !C$ या रूपातील सर्व
  !!{formula}s चे एक प्रगणन
  $!B_1 \lor !C_1$, $!B_2 \lor !C_2$, \dots,
  असू द्या. !!{formula}s च्या संचांची
  वाढती श्रेणी~$\Gamma_n$ परिभाषित करू;
  प्रत्येक $\Gamma_{n+1}$ हा $\Gamma_n$
  मध्ये एक नवे !!{formula} मिळवून
  परिभाषित केला जाईल.
  $\Gamma^*$ हा सर्व~$\Gamma_n$ चा
  संयोग असेल. नवी !!{formula}s अशी निवडू
  की $\Gamma^*$ अभाज्य राहील आणि
  $\Gamma^* \Proves/ !A$ सुद्धा राहील.
  यासाठी प्रत्येक पायरीवर पुढील अटी
  पूर्ण करणारा पहिला विकल्पयोग~
  $!B_i \lor !C_i$ शोधला पाहिजे:
  \begin{enumerate}
  \item\ollabel{gamma-1}
    $\Gamma_n \Proves !B_i \lor !C_i$
  \item\ollabel{gamma-2}
    $!B_i \notin \Gamma_n$ आणि
    $!C_i \notin \Gamma_n$
  \end{enumerate}
  $\Gamma_n \cup \{!B_i\} \Proves/ !A$
  असेल, तर $\Gamma_n$ मध्ये $!B_i$
  मिळवू; अन्यथा $!C_i$ मिळवू.
  हे यशस्वी होते ते दाखवायचे आहे.
  सध्या \olref{gamma-1} आणि~
  \olref{gamma-2} पूर्ण करणारा सर्वात
  लहान $i$ म्हणजे $i(n)$ असे ठरवू.

  $\Gamma_0 = \Gamma$ असे ठेवून
  \[
  \Gamma_{n+1} = \begin{cases}
    \Gamma_n \cup \{!B_{i(n)}\} &
    \text{जर $\Gamma_n \cup \{!B_{i(n)}\} \Proves/ !A$ असेल} \\
    \Gamma_n \cup \{!C_{i(n)}\} & \text{अन्यथा}
    \end{cases}
  \]
  अशी व्याख्या करू. $i(n)$ अपरिभाषित
  असेल, म्हणजे $\Gamma_n \Proves !B \lor !C$
  असताना नेहमी $!B \in \Gamma_n$
  किंवा $!C \in \Gamma_n$ असेल, तर
  $\Gamma_{n+1} = \Gamma_n$ ठेवू.
  आता $\Gamma^* = \bigcup_{n=0}^\infty \Gamma_n$
  असे ठेवू.

  प्रथम प्रत्येक $n$ साठी
  $\Gamma_n \Proves/ !A$ हे दाखवू.
  $n$ वर आगमन करू. $n = 0$ साठी
  हा दावा प्रमेयाच्या गृहीतकामुळे,
  म्हणजे $\Gamma \Proves/ !A$ मुळे,
  खरा आहे. $n>0$ असेल, तर
  $\Gamma_n \Proves/ !A$ असल्यास
  $\Gamma_{n+1} \Proves/ !A$
  हे दाखवायचे आहे.
  $i(n)$ अपरिभाषित असेल, तर
  $\Gamma_{n+1} = \Gamma_n$;
  काहीच सिद्ध करावे लागत नाही.
  म्हणून $i(n)$ परिभाषित आहे असे
  समजा. सोयीसाठी $i = i(n)$ असे लिहू.

  दाव्याचे व्यत्यस्त रूप सिद्ध करू.
  $\Gamma_{n+1} \Proves !A$ असे समजा.
  रचनेनुसार $\Gamma_n \cup
  \{!B_i\} \Proves/ !A$ असेल, तर
  $\Gamma_{n+1} = \Gamma_n \cup
  \{!B_i\}$; अन्यथा
  $\Gamma_{n+1} = \Gamma_n \cup \{!C_i\}$.
  पहिला प्रकार असूच शकत नाही;
  कारण तेव्हा $\Gamma_{n+1} \Proves/ !A$.
  म्हणून $\Gamma_n \cup
  \{!B_i\} \Proves !A$ आणि
  $\Gamma_{n+1} = \Gamma_n \cup \{!C_i\}$.
  $i(n)$ च्या व्याख्येनुसार
  $\Gamma _n \Proves !B_i \lor !C_i$.
  $\Gamma_n \cup \{!B_i\} \Proves !A$
  आहे. तसेच
  $\Gamma_{n+1} = \Gamma_n \cup
  \{!C_i\} \Proves !A$ आहे.
  त्यामुळे $\Gamma_n \Proves !A$,
  हेच दाखवायचे होते.

  $\Gamma^* \Proves !A$ असेल, तर
  $\Gamma' \subseteq \Gamma^*$
  असा काही परिमित उपसंच असेल की
  $\Gamma' \Proves !A$.
  प्रत्येक $!D \in \Gamma'$ कोणत्या
  तरी~$i$ साठी $\Gamma_i$ मध्ये असले
  पाहिजे. यांपैकी सर्वात मोठा
  $n$ घ्या. $i \le n$ असेल, तर
  $\Gamma_i \subseteq \Gamma_n$;
  म्हणून $\Gamma' \subseteq \Gamma_n$.
  पण मग $\Gamma_n \Proves !A$,
  आणि वर सिद्ध केलेल्या
  $\Gamma_n \Proves/ !A$ शी हा विरोधाभास.

  शेवटी $\Gamma^*$ अभाज्य आहे,
  म्हणजे \olref{defn:prime} मधील
  \olref{defn:prime1},
  \olref{defn:prime2} आणि~
  \olref{defn:prime3} या अटी
  पूर्ण करतो, हे दाखवू.

  प्रथम $\Gamma^* \Proves/ !A$;
  म्हणून $\Gamma^*$ सुसंगत आहे,
  आणि \olref{defn:prime1} पूर्ण होते.

  आता $\Gamma^* \Proves !B \lor !C$
  असेल, तर $!B \in \Gamma^*$
  किंवा $!C \in \Gamma^*$,
  हे दाखवू. यावरून
  \olref{defn:prime3} सिद्ध होते;
  कारण $!B \lor !C \in \Gamma^*$
  असेल, तर $\Gamma^* \Proves !B
  \lor !C$ सुद्धा. म्हणून
  $\Gamma^* \Proves !B \lor !C$
  पण $!B \notin \Gamma^*$
  आणि $!C \notin \Gamma^*$
  असे समजा.
  $\Gamma^* \Proves !B \lor !C$
  असल्याने कोणत्या तरी~$n$ साठी
  $\Gamma_n \Proves !B \lor !C$.
  $!B \lor !C$ हा विकल्पयोग
  प्रगणनात, समजा $!B_j \lor !C_j$
  म्हणून, येतो.
  $!B_j \lor !C_j$ हा
  $i(n)$ च्या व्याख्येतील गुणधर्म
  पूर्ण करतो: $\Gamma_n \Proves !B_j
  \lor !C_j$, तर
  $!B_j \notin \Gamma_n$ आणि
  $!C_j \notin \Gamma_n$.
  त्या निर्देशांकापेक्षा आधीचे निर्देशांक
  परिमित संख्येने आहेत. हा विकल्पयोग
  निवडला जात नाही तोवर अटी पूर्ण करत
  राहतो; तोवर निवडलेला विकल्पयोग
  $!B_i \lor !C_i$ त्याच्या आधीचाच असतो.
  त्यातील $!B_i$ किंवा $!C_i$
  मिळवल्यावर तो पुन्हा अटी पूर्ण करत नाही.
  म्हणून काही टप्प्यावर~$m$
  असे मिळेल की $j = i(m)$.
  पण तेव्हा $!B \in \Gamma_{m+1}$
  किंवा $!C \in \Gamma_{m+1}$;
  आणि हे $!B \notin \Gamma^*$
  व $!C \notin \Gamma^*$ या
  गृहीतकाच्या विरुद्ध आहे.

  आता $\Gamma^* \Proves !B$
  असे समजा. मग
  $\Gamma^* \Proves !B \lor !B$.
  पण $\Gamma^* \Proves !B \lor !B$
  असेल, तर $!B \in \Gamma^*$,
  हे आत्ताच सिद्ध केले आहे.
  त्यामुळे $\Gamma^*$ हा
  \olref{defn:prime} मधील
  \olref{defn:prime2} पूर्ण करतो.
\end{proof}

\begin{prob}
$\Gamma \Proves/ \bot$ असेल, तर
अभिजात तर्कशास्त्रात $\Gamma$ सुसंगत
असतो, म्हणजे~$\Gamma$ मधील सर्व सूत्रांना
सत्य करणारे मूल्यांकन अस्तित्वात असते,
हे दाखवा.
\end{prob}

\end{document}
